\manualchapter{Patch ideas and interaction}{sec:patches}
\subsection{Move a game variable with an LFO}
\begin{enumerate}
\item In Super Mario Bros., reach active play and SAVE a RackNES snapshot.
\item With the Genie row unplugged, select
\textbf{Player Horizontal Screen Position}.
\item Scale a slow LFO to 0--10 V, connect it, and watch the result.
\item Reduce the range to affect a smaller part of the byte's range.
\item Unplug the row before pressing RackNES LOAD to return to the saved
situation without immediately writing over it again.
\end{enumerate}
A memory write does not coordinate all of the game's other position or
collision variables. Jumps, wrapping, or unexpected movement can be a
consequence of editing one byte while the game is running.

\subsection{Explore Zelda's sound driver}
Select \textbf{The Legend of Zelda} with that game loaded in RackNES.
Start with one continuous entry, such as \textbf{Song Position (Pulse 1)}
or \textbf{Note Duration (Pulse 1)}, and a steady 0--10 V source. Use the
separate SQ1 output on RackNES to hear that voice in isolation.

These are addresses used by the game's sound driver, not general-purpose
synthesizer parameters. Not every byte represents a valid musical state.
The active song and game scene affect the result; names such as
\textbf{Frequency (Triangle)} do not imply V/oct tuning. Try a narrow range,
then compare with the input disconnected.

\subsection{Trigger a toggle}
Choose \textbf{Player Swimming} in Super Mario Bros. or a Zelda entry marked
T in the reference table. Feed separated 0--5 V gates. Each rising edge
alternates between the entry's two endpoint values. Holding a gate does not keep rewriting it,
and a second row's trigger state is independent.

\subsection{How rows interact}
If several rows target the same address, they compete. RackNES consumes rows
in top-to-bottom order, so a later row with a write in the same update takes
precedence. The game can then overwrite that value during emulation.
Assign distinct elements when you want independent control of game variables.

Genie works at Rack's sample rate rather than the controller inputs' slower
16-sample scan. Expander messages pass between processing steps, so a write
is not guaranteed to coincide with a particular NES instruction or frame.
RackNES processes Genie messages even during HANG. Use disconnected inputs,
not Hang, when you need to stop Genie from changing memory.
